% ---- Conclusion -------------------------------------------------------------
% Write last. Short (~0.5 page): restate contributions, key finding, future work.

\todo{Write after experiments are complete. Draft outline:}

We presented \method{}, a domain-specific language model for text-to-SVG icon
generation. We introduced the \dataset{} corpus --- \num{275912} license-filtered,
normalized SVG icons from the Iconify ecosystem, annotated with VLM-generated
natural language descriptions --- and demonstrated that LoRA fine-tuning of
Qwen3.5-9B on this corpus yields a model that generates valid, renderable,
semantically aligned SVG icons from natural language prompts.

\todo{Insert key numerical finding.}

\paragraph{Limitations.}
The \dataset{} corpus reflects the distribution of the Iconify ecosystem:
predominantly single-color, geometric icons in flat or outline styles. The
model may generalize poorly to photorealistic SVGs, complex illustrations,
or icon styles not represented in Iconify. Caption quality depends on the
vision model and prompt template; a small fraction of icons receive semantically
inaccurate descriptions that may reduce training signal quality.

\paragraph{Future work.}
We identify three promising directions. First, incorporating reasoning-enhanced
training \citep{svgen2025, reason_svg2025} via GRPO reinforcement learning may
further improve structural validity without increasing dataset size. Second,
interactive editing --- given an existing icon and a textual modification
instruction, produce an updated SVG --- is a natural extension enabled by the
instruction-following formulation. Third, evaluating the model in a
user study with designers and front-end developers would provide ecological
validity beyond the automated metrics used here.
